\vspace*{40mm}
\begin{center}
{\Large\bfseries ZUR\par\medskip HANNOVERSCHEN TRIANGULATION.}
\end{center}
\clearpage\null\clearpage
\begin{center}
[I.]\par\medskip
{\Large\bfseries BRIEFWECHSEL}\par\medskip
{\bfseries ZUR HANNOVERSCHEN TRIANGULATION.}\par\bigskip
[1.]\par\medskip
\textsc{Gauss} an \textsc{Schumacher}.\quad Göttingen, 5.~Juli 1816.
\end{center}

\par\noindent\hspace*{\parindent}Vor allem meinen herzlichen Glückwunsch zu dem herrlichen, großen Unternehmen, das Sie mir in Ihrem letzten Brief ankündigen. Diese Gradmessung in den königlich dänischen Staaten wird uns schon für sich genommen wertvolle Aufschlüsse über die Gestalt der Erde geben. Ich zweifle aber keineswegs daran, dass es in Zukunft möglich sein wird, Ihre Messungen durch das Königreich Hannover nach Süden fortzusetzen. Gegenwärtig kann ich einen solchen Wunsch in Hannover noch nicht zur Sprache bringen, da die Astronomie selbst zunächst noch so großer Unterstützung bedarf. Ich bin aber überzeugt, dass unsere Regierung, die ebenfalls gern die Wissenschaften unterstützt, bald dem ruhmreichen Beispiel Ihres vortrefflichen Königs folgen wird. Wir hätten dann bereits einen beträchtlichen Meridianbogen von \(6\frac{1}{2}\) Grad und könnten diese Arbeiten dann auch leicht mit den bayerischen Dreiecken verbinden. Letztere sind gewiss mit größter Sorgfalt gemessen worden, und es ist zu bedauern, dass sie der Öffentlichkeit vorenthalten werden.

\par\noindent\hspace*{\parindent}Für die Behandlung der gemessenen Dreiecke in der Rechnung habe ich eine eigene Methode entwickelt, die für einen Brief aber viel zu umfangreich wäre. Falls ich sie bis zu dem Zeitpunkt, zu dem Sie Ihre Dreiecke gemessen haben, noch nicht veröffentlicht haben sollte, werde ich sie ausführlich mit Ihnen besprechen. Ja, ich biete Ihnen an, die Berechnung der Hauptdreiecke selbst zu übernehmen.
\par\medskip\noindent IX.\hfill44
\clearpage
\noindent346\hfill{\small\textbf{BRIEFWECHSEL MIT SCHUMACHER.}}\par\medskip

\par\noindent\hspace*{\parindent}Beim zweiten Teil Ihres Unternehmens, der Messung des Längengrades, habe ich nur einen kleinen Zweifel. Ich meinte nämlich, dass die Länder der dänischen Monarchie eher \textls[120]{flach} zu nennen sind und jedenfalls keine hohen Berge haben. Wenn diese Annahme begründet ist und Sie deshalb zur Bestimmung des astronomischen Längenunterschieds einen Zwischenpunkt oder gar \textls[120]{mehrere} heranziehen müssen, so wird dieser Bestimmung immer eine kleine Unsicherheit anhaften, selbst wenn Sie noch so geschickte Gehilfen und noch so gute Hilfsmittel haben.

\par\noindent\hspace*{\parindent}Ein großer Vorteil für Sie ist, dass Dänemark schon einmal trigonometrisch vermessen worden ist. Ich meine damit natürlich nicht die gemessenen Winkel selbst, die für eine Gradmessung längst nicht ausreichen. Vielmehr werden Ihnen jene Arbeiten die Auswahl der Stationspunkte ungemein erleichtern. Diese Suche wäre für mich bei einer ähnlichen Arbeit gerade das Unangenehmste, weil dabei so viel Zeit umsonst verloren geht. Ich habe mir viel Mühe gegeben (mit Blick auf ähnliche künftige Arbeiten), die von \textsc{Epailly} im Hannoverschen gemessenen Winkel zu erhalten, aber ohne Erfolg.

\par\noindent\hspace*{\parindent}.\ .\ .\ .\ .\ .\quad Mir war eine interessante Aufgabe eingefallen, nämlich
\begin{quote}
\noindent»\textbf{allgemein} eine gegebene Fläche so auf eine andere (gegebene) zu projizieren (abzubilden), dass das Bild dem Original in den kleinsten Teilen ähnlich wird.«
\end{quote}
\par\noindent\hspace*{\parindent}Ein Sonderfall liegt vor, wenn die erste Fläche eine Kugel, die zweite eine Ebene ist. Hier sind die stereografische und die Mercator-Projektion partikuläre Lösungen. Gesucht ist aber für jede Art von Flächen die \textls[120]{allgemeine} Lösung, die alle partikulären Lösungen umfasst.

\par\noindent\hspace*{\parindent}Daran sollen \textsc{Monge} und \textsc{Poinsot} schon im Journal philomathique gearbeitet haben (wie \textsc{Burckhardt} an \textsc{Lindenau} geschrieben hat). Da ich aber nicht genau weiß, \textls[120]{wo}, habe ich noch nicht nachsehen können und weiß daher nicht, ob die Lösungen dieser Herren ganz meiner Idee entsprechen und die Sache erschöpfen. Andernfalls schiene mir dies einmal eine geeignete Preisfrage für eine Gesellschaft zu sein.\quad.\ .\ .\ .\ .\ .
